We proposed \tb, a structurally region-level cognitive map for active
spatial exploration, and motivated it by the two characteristic
failure modes the recent \tos benchmark identifies in foundation
models---belief drift and belief inertia. We formalised the
representation, derived analytical predictions on \tos's Stability
and Belief~Inertia metrics
(Propositions~\ref{prop:stab}--\ref{prop:inertia}), implemented the
representation as a post-hoc replay tool that ingests \tos
exploration logs without retraining, and validated the pipeline on a
synthetic trajectory. Full-scale runs against Qwen-VL-7B on 25
procedural 3-room scenes (text and vision worlds) are scheduled for
the next two days.

\paragraph{Looking ahead.}
Three concrete extensions follow from this proposal:
(i) \emph{familiarity-weighted overwrite}, in which territory
familiarity tightens the contradiction tolerance (sketched in
Section~\ref{app:math});
(ii) \emph{multi-agent territorial belief}, in which two agents share
territory ownership and revise each others' beliefs through
familiarity-aware handoff (a direction \tos itself flags as
future work);
(iii) \emph{learned-not-just-replayed} \tb, in which the
representation is wired into a VAGEN \citep{wang2025vagen} world
model so that the agent itself \emph{acts} territorially, rather than
the LLM acting flatly while \tb merely observes. Each of these
extends the diagnostic baseline of this proposal into a substantive
research line; the present project's contribution is to establish
that the diagnostic question is well-posed, the formal model exists,
and the implementation has been validated end-to-end.
